\documentclass[12pt]{article}
\usepackage[utf8]{inputenc}
\usepackage[T1]{fontenc}
\usepackage{hyperref}
\title{Borgu Fulfulde Language Resources for K-12 Curriculum Development}
\author{Compiled by Claude Sonnet 5.5 (claude-sonnet-5-5) for Andr\'as Kornai}
\date{}
\begin{document}
\maketitle

\section{Introduction}
Borgu Fulfulde (ISO 639-3: fue; Glottocode: borg1235) is a Fula variety spoken mainly in the Borgou Department of Benin, with speakers also in Nigeria, Togo and Burkina Faso. Dialects listed are Bakuure, Korakuure, Juguure and Caabankeere. It is written in Latin script. The search found mostly Christian mission materials and no secular educational resources.

\section{Primary Educational Resources}
None found in the search results (no primers, textbooks or school programmes).

\section{Linguistic and Reference Materials}
\begin{itemize}
\item Wikipedia article on Borgu Fulfulde (background, dialects): \url{https://en.wikipedia.org/wiki/Borgu_Fulfulde}
\end{itemize}
No dictionary or grammar was found.

\section{Religious and Community Resources}
\begin{itemize}
\item SIM story on the completed Borgu Fulfulde Bible (New Testament 2011, Old Testament completed early 2024): \url{https://www.sim.org/stories/borgu-fulfulde-bible-translation-completed-benin-fulani/}
\item Scripture Earth listing for fue: \url{https://www.scriptureearth.org/00eng.php?iso=fue}
\end{itemize}

\section{Digital and Audio Resources}
\begin{itemize}
\item Global Recordings Network page, Fulfulde, Borgu: \url{https://globalrecordings.net/en/language/fue}
\item Global Recordings Network page, Fulfulde, Borgu: Bakuure (no recordings listed): \url{https://globalrecordings.net/en/language/9922}
\end{itemize}

\section{Curriculum Development Context}
The complete Bible is the largest text corpus found and could supply reading material, with the orthography it uses as a starting point. Secular and school materials appear to be absent from what was found and would need to be created.

\end{document}